% TRANSLATION UNIT: translate this file only.
% ===========================================================================
%  02 — Justificación
%  The \section title is the canonical (Spanish) chapter name; the body prose
%  is written in ENGLISH and translated to Spanish later.
% ===========================================================================

% DONE: justification written in English, four strands plus an explicit section
% of disclaimers. Sources: banrep2024financial (abstract only, figure and date
% read on the institution's page) and lusardi2019financial (the tailoring
% conclusion only, already reviewed in chapter 04). No claim rests on a source
% whose body was not read.
%
% CLAIMS INVENTORY — every prose claim below with its evidence.
% EVIDENCE: 16.4 % answer the Big Three correctly → banrep2024financial abstract,
%   page read 2026-10-27; figure and publication date (2024-02-29) visible on the
%   page; body of the working paper NOT retrievable (refs/fuentes-acceso.md:85).
% EVIDENCE: ten seed definitions, four Colombian → services/simulator/src/domain/
%   seeds/mod.rs:145-155 (drafts()); colombia.rs:1-3,212,302.
% EVIDENCE: 24 % EA equals 1.809 % monthly, not 2 % → seeds/colombia.rs:17-19.
% EVIDENCE: GMF exemption 350 UVT read from @UVT; rate 0.004 as a literal of the
%   law (Estatuto Tributario art. 872) → seeds/colombia.rs:44-56,212-230.
% EVIDENCE: CDT agreed in days, quoted in EA, simple interest overstates the
%   return → seeds/colombia.rs:23-25,57-62,280-283.
% EVIDENCE: native code took the UVT as a per-call parameter, so the calculator
%   was silently wrong every January 1st → seeds/colombia.rs:31-38.
% EVIDENCE: financial_indicators: NUMERIC(20,6), DATERANGE, EXCLUDE USING gist
%   → services/simulator/migrations/20260902121500_financial_indicators.up.sql:1-44.
% EVIDENCE: calculators are data; version published separately → migrations
%   20260902120000_calculators.up.sql, 20260902130000_calculator_published_version.up.sql.
% EVIDENCE: builder RPCs (Upsert/Validate/Submit/Approve/Reject) →
%   contracts/proto/fintcart/simulator/v1/simulator.proto:32-54.
% EVIDENCE: categories catalog and versioned body_doc → services/learning/
%   migrations/20260902100000_add_categories.up.sql:13-34.
% EVIDENCE: saga compensations exist and are tested → services/orchestrator/
%   internal/server/saga.go, saga_anonymization_test.go, saga_grading_test.go,
%   saga_registration_test.go.
% EVIDENCE: f32/f64 prohibited as a hard error in CI → .github/workflows/ci.yml
%   (clippy), reported in chapter 11.
% EVIDENCE: seed regression pins seeds to the native oracle → tests/seed_regression.rs.
% EVIDENCE: four developer verbs, zero manual steps → dev/build:20-26, dev/up:15-38,
%   dev/migrate, dev/down; CI verifies the three verbs still run untouched.
% EVIDENCE: eight services, none of which bills → dev/docker-compose.yaml; services/.
% EVIDENCE: layout tested down to 360 px → frontend/e2e/visual/*.spec.ts WIDTHS.

\section{Justificación}

This chapter states why the platform is worth building, for whom, and on what
evidence. Chapter~04 establishes that financial literacy is low where measured, that
education moves knowledge more reliably than behavior, and that digital delivery
has firmer support than gamification, so none of that is reargued here. What
remains is narrower: that this artifact, built here and now, lets a reader or an
institution do something they could not otherwise do. The proposal set out four
strands---social relevance, practical implications, theoretical value, and
convenience---and they are taken in that order. Each claim is tied to something
that exists in the repository rather than to something planned, and the chapter
closes by listing what the work does not claim, because the boundary is part of
the argument.

\subsection{Social relevance}

The Banco de la República published in February 2024 a working paper on
financial literacy and financial wellbeing in Colombia under high inflation. Its
abstract reports that 16.4~\% of Colombians answer the three questions known as
the Big Three---compound interest, inflation, and risk diversification---all
correctly \parencite{banrep2024financial}. The figure and the date are on the
institution's page; the paper's body is not, because its PDF is served through a
script the archive did not retrieve. This chapter therefore cites the number
and not the study's sampling design, and says so rather than implying a reading
it did not do. Sixteen in a hundred is not a marginal deficit: for five of every
six people the arithmetic behind a loan's cost or a deposit's return is
unavailable, so the decision is made on the price tag instead of the rate.

That deficit is where the project's first design choice comes from. A foreign
personal-finance application is built around its own country's instruments: a
tax-deferred retirement account, an escrowed-property-tax mortgage, a national
credit score. A Colombian reader holds none of those. What they hold is a
consumer loan whose rate the Superintendencia Financiera requires the lender to
publish as an \emph{efectiva anual} while the instalment is settled on a nominal
periodic rate; an account charged four pesos per thousand on every withdrawal,
with a monthly exemption measured in UVT; a UVT that changes every January; and
a CDT agreed in days but quoted in effective annual terms. The simulator ships
ten seed definitions, four of them Colombian
(\code{services/simulator/src/domain/seeds/colombia.rs});
Table~\ref{tab:semillas-co} states what each computes and which misreading it
prevents.

\begin{table}[htbp]
  \centering
  \caption{The four Colombian seed definitions and the misreading each one prevents.}
  \label{tab:semillas-co}
  \begin{tabular}{@{}p{2.5cm}p{5.4cm}p{8.5cm}@{}}
    \hline
    Seed & What it computes & The error it prevents \\
    \hline
    \code{ea_a_mv} & The nominal periodic rate implied by an effective annual
    rate. & Dividing a 24~\% E.A. by twelve and quoting 2~\% monthly, when
    the equivalent monthly rate is 1.809~\%. \\
    \code{mv_a_ea} & The effective annual rate implied by a nominal periodic
    rate, so two offers can be compared. & Treating a monthly quote and an
    annual one as the same unit. \\
    \code{gmf} & The four-per-thousand levy on a movement, net of the 350 UVT
    monthly exemption, with the UVT read from the \code{@UVT} indicator. & A
    balance projected as if withdrawals were free, and an exemption computed
    with last year's UVT. \\
    \code{cdt_dias} & The return on a CDT agreed for a term in days, at an
    effective annual rate, with withholding at source. & \code{capital * tasa *
    dias / 365}, which is simple interest and overstates the yield of a long
    term. \\
    \hline
  \end{tabular}
\end{table}

Read as a whole, the table makes the argument that translation cannot: a
localized interface translates labels, not arithmetic. No foreign tool computes
the four-per-thousand levy, because no foreign tax authority imposes it, and
none measures an exemption in a unit republished annually. The rate-convention
pair is the sharper case: one institution describes the same loan in two units,
and the conversion is the concept the student is missing. This is the design
lesson that programs must be tailored to an identified audience
\parencite{lusardi2019financial}, applied to the instrument rather than to the
message. The platform also charges nothing, and the topology is the evidence:
eight services, none of which bills or advertises. Access is a browser and a
Spanish interface, and the layout is tested down to 360~px, so it fits the phone
a student carries. That is a claim about reach, not outcomes.

\subsection{Practical implications}

The second strand is what an institution can do with the artifact; each
capability replaces a different bottleneck. First, publish curated content
without a release. Feature 002 replaced the free-text category and the plain-text
body with a managed category catalog and a versioned block document
(\code{services/learning/migrations/20260902100000_add_categories.up.sql}), with
editors drafting and an editorial coordinator publishing, so an article is
reviewed and published by the people responsible for it and no release is needed
to correct a paragraph.

Second, add a calculator without a deploy. A calculator is data: an identity row
plus an immutable definition version, created on the builder surface, validated
before it is stored, submitted for review, and approved by the editorial
coordinator (\code{simulator.proto:32-54}). A new formula enters the catalog as
a row that passed the same validation as every other row, not as a release that
passed a code review. The seed definitions are written in that language and
parsed by that validator, so the built-in calculators are not exempt from the
checks they impose on users.

Third, keep annual indicators current with validity ranges. The UVT and the
minimum wage are republished every year, and the earlier implementation took the
UVT as a parameter supplied by the caller, which left the calculator silently
wrong from January 1st onward when the caller had not updated it. Indicators now
live in \code{financial_indicators} with \code{NUMERIC(20,6)} values and a
\code{DATERANGE} validity, and a database exclusion constraint forbids
overlapping periods for the same indicator
(\code{20260902121500_financial_indicators.up.sql:44}). Because each stored
simulation keeps the indicator values it used, a run from 2026 can still be
explained in 2029: definition version and values are recoverable. For an
institution, that is the difference between a number it published and a number
it can defend.

\subsection{Theoretical value}

The theoretical contribution is small, and inflating it would be the easiest
dishonesty available here. The project proposes no learning theory, does not
claim that microservices teach better, and adds nothing to the evidence that
gamification motivates. What it demonstrates is narrower: that the constraints
normally justified by banking regulation---exact decimal arithmetic, a database
per service, compensation instead of distributed transactions---carry into an
educational product without the product becoming a bank, and that those
constraints are what make its numbers trustworthy rather than what they cost.

Monetary values are \code{rust_decimal} at every layer, and the continuous
integration job that runs \code{clippy} turns the use of \code{f32} or
\code{f64} into a hard error rather than a review comment. Data ownership is a
schema boundary, not a convention: the simulator knows nothing about articles and
the learning service knows nothing about rates. Operations
that span services are compensated rather than assumed atomic, and the
compensations have their own tests for registration, grading, and anonymization
(\code{services/orchestrator/internal/server/}). The payoff is not decorative: a
student who sees a rounded figure learns a rule that will be wrong at the
counter, and a seed regression suite exists so that no future edit can move a
published result silently (\code{services/simulator/tests/seed_regression.rs}).
Engineering rigor here is ordinary discipline applied where the audience cannot
check the arithmetic themselves.

\subsection{Convenience and practicality}

The platform is reproducible from four verbs and a container runtime.
\code{dev/build} builds the images and prints the next command, \code{dev/up}
raises the topology and waits on health checks, \code{dev/migrate} applies every
migration, and \code{dev/down} removes it. No manual step separates them, and
that claim is itself tested: a continuous integration job runs the verbs to
confirm the documented workflow has not drifted from the scripts. A reader who
wants to check a figure in this report can raise the same stack, run the suite
that produced it, and disagree with the result.

The software is also measurable from outside. The verification suites, the load
profile, the visual probes at five widths, and the accessibility assertions are
described in chapter~11 with enough detail to be re-run; the human-subject
instruments---consent, a perception survey, a moderated session with the System
Usability Scale---exist as a closed package rather than a promise. A reader
therefore does not have to trust the author's impression of a screen, and the
author does not have to defend an impression a tool could have settled.

\subsection{What this work does not claim}

The last movement is subtraction. This platform does not replace formal
financial education, and nothing in the previous chapters suggests a browser
application can do what a curriculum, a teacher, and time do. It does not advise
anyone on their money: it recommends no product, ranks no option, and tells no
student what to do with their salary. A platform that simulates is not a
platform that recommends. It computes the consequences of assumptions the
student supplies, shows the arithmetic, and stops, because the assumptions
belong to the person and responsibility for the decision is not transferred
with the result. Nor does it claim to move behavior: chapter~04 reported that
financial education improves knowledge more reliably than behavior, and this
chapter does not upgrade that finding into a promise about habits. It makes no
claim on national literacy rates, and none that a student who uses it will be
better off than one who does not. It offers ten calculators, four of them
locally calibrated, a way to publish and update them without a release, and a
stack a reader can rebuild and re-measure. That is a modest thing to defend, and
it is what the evidence supports.
